\chapter{Hypothermia}

\section*{Definition}
Hypothermia is a core body temperature below 35 C (95 F). It is classified as mild (32-35 C), moderate (28-32 C), or severe (below 28 C). Hypothermia is a life-threatening emergency that requires immediate intervention.

\section*{What Happens in Your Body}
When the body loses heat faster than it can produce it, thermoregulatory mechanisms are overwhelmed. Initial shivering generates heat but becomes ineffective as temperature drops. Below 32 C, shivering ceases and metabolic rate declines by about 5 percent per degree. Cardiac conduction becomes unstable, with risk of ventricular fibrillation. Coagulation is impaired. Cerebral blood flow decreases, causing progressive loss of consciousness.

\section*{Causes}
\begin{itemize}
\item Prolonged exposure to cold environments
\item Cold water immersion (heat loss is 25x faster in water)
\item Inadequate clothing in cold weather
\item Homelessness or inadequate shelter
\item Advanced age (decreased thermoregulation)
\item Alcohol or drug intoxication (widening of blood vessels, impaired judgment)
\item Hypoglycemia
\item Hypothyroidism or adrenal insufficiency
\item Sepsis
\item Burns or extensive skin disease
\item Spinal cord injury or immobility
\end{itemize}

\section*{When to See a Doctor}
\begin{itemize}
\item Shivering (mild) or no shivering (severe)
\item Cold skin to the touch
\item Confusion, apathy, or drowsiness
\item Slurred speech and poor coordination
\item Weak pulse and slow, shallow breathing
\item Hypotension
\item Arrhythmias (atrial fibrillation, ventricular fibrillation)
\item Coma or unresponsiveness in severe hypothermia
\end{itemize}

\section*{Related Symptoms}
\begin{itemize}
\item Frostbite (local cold injury)
\item Chilblains (pernio)
\item Trench foot
\item Raynaud phenomenon
\item Bradycardia
\item Hypoglycemia
\item Alcohol intoxication
\end{itemize}

\section*{Self-Care/Home Management}
\begin{itemize}
\item Call emergency services - hypothermia is a medical emergency
\item Move the person to a warm environment
\item Remove wet clothing and replace with dry, warm layers
\item Apply warm compresses to the trunk (not extremities)
\item Provide warm, non-alcoholic, non-caffeinated drinks if conscious
\item Do NOT apply direct heat (heating pads, hot water)
\item Do NOT rub or massage limbs
\item Handle the person gently to avoid triggering arrhythmias
\end{itemize}

\section*{How Your Doctor Will Evaluate You}
\begin{itemize}
\item Core body temperature measurement with esophageal or rectal thermometer
\item Clinical evaluation of severity based on mental status and shivering
\item Electrocardiogram (Osborn J wave is characteristic)
\item Complete blood count and electrolytes
\item Coagulation profile
\item Blood glucose
\item Arterial blood gas analysis
\item Renal and liver function tests
\end{itemize}

\section*{Treatment Options}
\begin{itemize}
\item Passive external rewarming (warm environment, blankets) for mild hypothermia
\item Active external rewarming (forced warm air, heating pads) for moderate hypothermia
\item Active internal rewarming (warmed IV fluids, warmed humidified oxygen, peritoneal lavage) for severe hypothermia
\item Extracorporeal membrane oxygenation (ECMO) for refractory severe hypothermia
\item Cardiac monitoring and careful handling to prevent arrhythmias
\item Treat underlying cause (hypoglycemia, hypothyroidism)
\item Avoid rewarming shock: slow, controlled rewarming
\end{itemize}
